\chapter{Esophageal Spasm}

\section*{Definition}
Esophageal spasm encompasses hypercontractile motility disorders of the esophagus, including diffuse esophageal spasm (DES) and jackhammer esophagus, characterized by simultaneous or high-amplitude contractions. Diagnosis requires high-resolution manometry showing premature contractions (distal latency less than 4.5 s) with 20\% or greater of swallows.

\section*{What Happens in Your Body}
Disordered inhibitory neural signaling (nitric oxide pathway) in the esophageal myenteric plexus results in loss of peristaltic coordination and premature, simultaneous, or repetitive contractions. Increased cholinergic excitation may contribute to the high-amplitude contractions seen in jackhammer esophagus.

\section*{Causes}
\begin{itemize}
  \item \textbf{Idiopathic:} Most common cause; no identifiable underlying disorder
  \item \textbf{GERD:} Chronic acid reflux may trigger secondary esophageal spasm through sensitization of esophageal mechanoreceptors
  \item \textbf{Neuropathy:} Diabetic autonomic neuropathy, Chagas disease (\textit{Trypanosoma cruzi}) damaging the myenteric plexus
  \item \textbf{Obstructive:} Mechanical obstruction (stricture, ring, tumor) causing secondary hypercontractility
  \item \textbf{Other:} Anxiety, panic disorder, visceral hypersensitivity, eosinophilic esophagitis, paraneoplastic syndrome
\end{itemize}

\section*{When to See a Doctor}
See your doctor if:
\begin{itemize}
  \item Retrosternal chest pain mimicking cardiac angina (must rule out myocardial ischemia first)
  \item Dysphagia to both solids and liquids (intermittent or progressive)
  \item Odynophagia or regurgitation of undigested food
  \item Chest pain with swallowing or that awakens you from sleep
  \item Failure to respond to empiric acid suppression therapy
\end{itemize}

\section*{Related Symptoms}
\begin{itemize}
  \item Retrosternal chest pain, dysphagia, odynophagia, globus sensation, regurgitation, heartburn
\end{itemize}
\textbf{Important:} Esophageal spasm can present identically to cardiac chest pain; all people with undifferentiated chest pain should have ECG and cardiac biomarkers before esophageal evaluation.

\section*{Self-Care / Home Management}
\begin{itemize}
  \item Avoid trigger foods: very cold or very hot liquids, carbonated beverages, acidic foods, red wine
  \item Eat slowly, chew thoroughly, and take small bites
  \item Upright posture during and for 30 minutes after meals
  \item Stress reduction techniques (deep breathing, meditation) as anxiety often exacerbates symptoms
  \item Keep a symptom diary to identify individual triggers
\end{itemize}

\section*{How Your Doctor Will Evaluate You}
\begin{itemize}
  \item High-resolution esophageal manometry (HRM) with Chicago Classification: DES shows premature contractions (distal latency less than 4.5 s) with 20\% or greater of swallows; jackhammer shows a distal contractile integral (DCI) greater than 8000 mmHg{$\cdot$}s{$\cdot$}cm
  \rule{0pt}{0pt}
  \item EGD with biopsy to exclude mechanical obstruction, eosinophilic esophagitis, or malignancy
  \item Barium esophagram may show corkscrew esophagus in DES (simultaneous contractions)
  \item Ambulatory 24-hour pH monitoring to evaluate for GERD as a contributing factor
  \item Cardiac workup (ECG, echocardiogram, stress test) before attributing chest pain to esophageal spasm
\end{itemize}

\section*{Treatment Options}
\begin{itemize}
  \item Calcium channel blockers (nifedipine 10--30 mg or diltiazem 60--120 mg before meals) to reduce contraction amplitude
  \item Nitrates (isosorbide dinitrate 5--10 mg sublingual before meals) for acute symptom relief
  \item Proton pump inhibitors for concurrent GERD, which may trigger or exacerbate spasm
  \item Low-dose tricyclic antidepressants (imipramine 25--50 mg, trazodone) for visceral pain modulation
  \item Peroral endoscopic myotomy (POEM) for severe, medically refractory DES or jackhammer esophagus
\end{itemize}

\section*{References}
\begin{thebibliography}{99}
\bibitem{ref1} Pandolfino JE, Gawron AJ, Achkar E, et al. "Esophageal Spasm: Diagnosis and Management." Am J Gastroenterol. 2014;109(9):1336--1348.
\bibitem{ref2} Roman S, Pandolfino JE. "Distal Esophageal Spasm." In: Castell DO, Richter JE, eds. The Esophagus. 5th ed. Wiley-Blackwell; 2012.
\bibitem{ref3} Spechler SJ, Castell DO. "Classification of Oesophageal Motility Abnormalities." Gut. 2001;49(1):145--151.
\bibitem{ref4} Kahrilas PJ, Bredenoord AJ, Fox M, et al. "The Chicago Classification of Esophageal Motility Disorders." Neurogastroenterol Motil. 2015;27(2):150--162.
\bibitem{ref5} Richter JE. "Oesophageal Motility Disorders." Lancet. 2001;358(9284):823--828.
\end{thebibliography}
\clearpage
